% app_shell --- reference manual
%
% Build:  pdflatex app_shell.tex   (twice, for the table of contents)
% To Markdown:  pandoc app_shell.tex -o app_shell.md --standalone
%
% Deliberately plain LaTeX: article class, no custom macros beyond one listing
% style, so pandoc converts it without losing structure.

\documentclass[11pt,a4paper]{article}

\usepackage[utf8]{inputenc}
\usepackage[T1]{fontenc}
\usepackage[margin=2.6cm]{geometry}
\usepackage{booktabs}
\usepackage{longtable}
\usepackage{xcolor}
\usepackage{listings}
\usepackage{enumitem}
\usepackage[hidelinks]{hyperref}

\definecolor{codebg}{gray}{0.96}
\lstset{
  basicstyle=\ttfamily\small,
  backgroundcolor=\color{codebg},
  frame=none,
  breaklines=true,
  columns=fullflexible,
  keepspaces=true,
  showstringspaces=false,
  language=C++,
  captionpos=b,
  xleftmargin=1em,
}

\title{\texttt{app\_shell}\\[0.3em]\large The application shell:\\
one program, three platforms}
\author{Reference manual}
\date{}

\begin{document}
\maketitle
\tableofcontents
\newpage

%======================================================================
\section{What this library is}

\texttt{app\_shell} is the layer that turns ``I can draw'' into ``I have a
program''.

It sits in a family of three first-party libraries with a strict division of
labour:

\begin{center}
\begin{tabular}{@{}lll@{}}
\toprule
Library & Answers & Gives you \\
\midrule
\texttt{audio\_engine} & how does it sound?   & decoders, DSP, USB/ALSA/JACK/WASAPI/AAudio \\
\texttt{vk\_canvas}    & how does it look?    & Renderer, Canvas, text, Vulkan on three platforms \\
\texttt{app\_shell}    & how does it \emph{run}? & window, message pump, input, paths, scale \\
\bottomrule
\end{tabular}
\end{center}

The distinction between the last two is the point of this document.
\texttt{vk\_canvas} alone gives you a picture: a surface, a renderer, and
primitives to draw with. It does not give you a program. Somebody still has to
create the real window, write the Win32 message loop, write the Wayland event
loop, write Android's \texttt{ALooper} loop, install a crash handler, open a log
file, decide where files live, and decide how big the text should be on a screen
nobody predicted. That is roughly 1\,400 lines of platform code, and it is what
this library contains.

\subsection{The problem it exists to solve}

Before it existed, an application that wanted desktop \emph{and} Android had to
copy the previous application's \texttt{gui/} directory. That is not reuse. By
the third copy the three have diverged, and a fix in one does not reach the
others.

The specific failure it prevents is worse than duplication, though. The three
platform backends in the original code each dispatched into the application by
calling its methods \emph{by name}. Measured across them, they called nineteen
methods any application would have (\emph{the pointer moved}, \emph{a key went
down}, \emph{the surface died}) and seven that only a music player could. Five
of those seven were the same \texttt{switch} statement written out four times
over --- once per backend, plus once more in a headless screenshot tool. They
could not drift, being identical; they could only all have to change together,
which is worse.

\subsection{What it is not}

It is not a UI toolkit. It has no widgets, no layout engine, no theme. It draws
nothing. Those belong to \texttt{vk\_canvas} and to the application.

It is not an application framework in the inversion-of-control sense: it does
not own \texttt{main}'s loop and call you back. It owns the platform
\emph{bootstrap} (\texttt{main}/\texttt{WinMain}) and then hands control to the
application, which writes its own frame loop over \texttt{Host::pump()}. See
Section~\ref{sec:loop}.

%======================================================================
\section{Getting started}

\subsection{The smallest application}

Two files. First, the application itself:

\begin{lstlisting}
#include "app_main.hh"   // app_shell_main(), which YOU define
#include "app_view.hh"   // the interface the platform calls
#include "host.hh"       // the interface you call

class MyApp : public AppView {
public:
    bool create() {
        host_ = make_host();          // the platform's real Host
        if (!host_) return false;
        if (!host_->init(this)) return false;
        host_->showWindow();
        return true;
    }

    void run() {
        while (running_) {
            host_->pump(/*haveWork=*/dirty_);
            if (host_->quitRequested()) break;
            if (dirty_) { draw(); dirty_ = false; }
        }
    }

    // --- AppView ---
    void onHostResized() override { dirty_ = true; /* relayout */ }
    void shutdown()      override { running_ = false; }
    void onLButtonDown(int x, int y) override { /* ... */ dirty_ = true; }

private:
    void draw() { /* vk_canvas: Canvas -> Renderer::draw() */ }

    std::unique_ptr<Host> host_;
    bool running_ = true;
    bool dirty_   = true;
};

int app_shell_main() {
    MyApp app;
    if (!app.create()) return 1;
    app.run();
    return 0;
}
\end{lstlisting}

Note what is \emph{absent}: there is no \texttt{main()}, no \texttt{WinMain()},
no \texttt{\#ifdef \_WIN32}, and no mention of Wayland, Win32 or Android
anywhere. \texttt{app\_shell} supplies the entry point for both desktops and
calls \texttt{app\_shell\_main()} once the platform is ready.

\subsection{And the CMake}

\begin{lstlisting}[language=]
set(APP_SHELL_APP_NAME  "my_app")        # names my_app.log
set(APP_SHELL_WIN_CLASS "MyAppMain")     # the Win32 window class
add_subdirectory(path/to/app_shell app_shell_build)

add_executable(my_app src/my_app.cc)
target_link_libraries(my_app PRIVATE
    app_shell
    $<IF:$<PLATFORM_ID:Windows>,app_shell_win32,app_shell_wayland>
    vk_canvas_core)
\end{lstlisting}

That is the whole desktop story. Android costs one more file, described in
Section~\ref{sec:android}, and it is six lines long.

%======================================================================
\section{Architecture}

\subsection{The seam has two halves}

\begin{center}
\begin{tabular}{@{}lll@{}}
\toprule
Header & Implemented by & Called by \\
\midrule
\texttt{host.hh}     & the platform     & the application \\
\texttt{app\_view.hh} & the application & the platform \\
\bottomrule
\end{tabular}
\end{center}

That is the entire rule, and it is worth stating so plainly because getting the
direction of a call wrong is how platform code ends up knowing what an album is.

\texttt{Host} is what the application \emph{asks for}: give me a window, start a
timer, change the cursor, keep the display awake, wake me from another thread.
\texttt{AppView} is what the platform \emph{reports}: the pointer moved, a key
went down, you were resized, your drawing surface just died.

\subsection{Ownership and lifetime}

The application owns the \texttt{Host} (typically a
\texttt{std::unique\_ptr<Host>} member). The \texttt{Host} holds a non-owning
\texttt{AppView*} given to it by \texttt{init()}. There is exactly one of each
per process on every supported platform.

\texttt{make\_host()} is a free function each backend defines. It returns the
platform's real host on Windows and Linux; \textbf{on Android it returns
\texttt{nullptr} on purpose}, because an \texttt{AndroidHost} needs the
\texttt{android\_app*} that only \texttt{android\_main()} is ever handed. There
the host is \emph{injected} instead. A caller that forgets fails at startup with
a log line rather than dereferencing null.

\subsection{The frame loop}
\label{sec:loop}

\texttt{app\_shell} does not provide one. The application writes it, because
what counts as ``work to do'' is the application's question --- an editor idles
at zero frames, a game never idles.

The pattern every current consumer uses is render-on-demand:

\begin{lstlisting}
while (running_) {
    bool haveWork = renderer_ && pendingFrames_ > 0;
    host_->pump(haveWork);              // blocks when there is nothing to do
    if (host_->quitRequested()) break;
    if (!renderer_) continue;           // re-tested AFTER the pump; see below
    if (pendingFrames_ > 0) { drawFrame(); --pendingFrames_; }
}
\end{lstlisting}

Passing \texttt{haveWork == false} lets \texttt{pump()} block instead of
busy-spinning, so an idle application costs approximately zero CPU. Timers and
cross-thread events wake the wait normally.

\subsection{Two rules learned by crashing on a phone}

These are not style advice. Each corresponds to a real crash.

\begin{enumerate}[leftmargin=*]
\item \textbf{Nothing decided before \texttt{pump()} may be trusted after it.}
  \texttt{pump()} is where a platform delivers \emph{your surface is gone}. A
  \texttt{bool canDraw} computed before the pump and used after it segfaults on
  the first launch that hides the system bars, because hiding them recreates the
  window. Re-test after the pump, always.

\item \textbf{\texttt{safeInsets()} is the display cutout and never a system
  bar.} A bar is software: hide it and the screen is genuinely yours, and an
  inset for a bar that is not on screen leaves a permanent empty strip. A cutout
  is glass --- it is still there when every bar is hidden, and nothing drawn
  under it can be read.
\end{enumerate}

\subsection{Three vocabularies the host carries and never reads}

The mechanism that keeps a message pump from learning an application's domain:
three identifiers are passed \emph{through} the platform as plain integers.

\begin{description}[leftmargin=1.5em]
\item[App events] \texttt{postAppEvent(id, p1, p2)} comes back out of
  \texttt{onAppEvent(id, p1, p2)}, unchanged, on the UI thread. No backend
  branches on \texttt{id}.
\item[Timer ids] \texttt{startTimer(id, ms)} comes back as
  \texttt{onTimer(id)}.
\item[Hotkey ids] \texttt{registerHotkey(id, keyCode)} comes back as
  \texttt{onHotkey(id)}.
\end{description}

All three were once enumerations declared inside \texttt{host.hh}, with names
like \texttt{TrackChange} and \texttt{SeekUpdate}, and every backend switched on
them.

The same principle covers \texttt{launchArgument()}: a platform may be started
already knowing something (Android carries extras on its launch intent). The
host \emph{states the string} and stops. What it \emph{means} is decided in
\texttt{AppView::onHostReady()}.

\begin{quote}
\textbf{The test for anything added to this library:} could a drawing program
use it? If the answer needs a word from one specific application's domain, it
belongs in that application.
\end{quote}

%======================================================================
\section{Reference: \texttt{Host}}

Everything the application asks of the platform. Pure virtual unless a default
is shown; a defaulted method is one an honest platform may genuinely have no
answer for.

\subsection{Lifecycle}

\begin{longtable}{@{}p{0.42\textwidth}p{0.52\textwidth}@{}}
\toprule
\textbf{Method} & \textbf{Contract} \\
\midrule
\endhead
\texttt{bool init(AppView* owner)} &
Creates the native window and binds the owner. Returns false on failure.
Everything \texttt{pump()} later dispatches goes to this owner. \\
\addlinespace
\texttt{void showWindow()} &
Make the window visible. Separate from \texttt{init()} so an application can
finish building before anything is on screen. \\
\addlinespace
\texttt{void pump(bool haveWork)} &
One iteration of the platform's message loop. Blocks when \texttt{haveWork} is
false; returns promptly otherwise. Dispatches into the owner. \\
\addlinespace
\texttt{bool quitRequested() const} &
True once the platform has asked the process to end. \\
\bottomrule
\end{longtable}

\subsection{Assets and paths}

\begin{longtable}{@{}p{0.42\textwidth}p{0.52\textwidth}@{}}
\toprule
\textbf{Method} & \textbf{Contract} \\
\midrule
\endhead
\texttt{std::string exeDir() const} &
Directory holding the executable, UTF-8, \emph{with} a trailing separator.
\textbf{Empty on Android}, which is deliberate --- see below. \\
\addlinespace
\texttt{SurfaceProvider\& surfaceProvider()} &
\texttt{vk\_canvas}'s seam: instance extensions, \texttt{VkSurfaceKHR}, extent. \\
\addlinespace
\texttt{AssetReader\& assetReader()} &
\texttt{vk\_canvas}'s shipped data, rooted at \texttt{<exe>/assets/} --- shaders. \\
\addlinespace
\texttt{AssetReader\& dataReader()} &
The \emph{application's} read-only shipped data (fonts), at paths relative to
\texttt{exeDir()}. \\
\bottomrule
\end{longtable}

\paragraph{Why \texttt{dataReader()} is separate, and why \texttt{exeDir()} is
empty on Android.} On both desktops these are plain filesystem reads and the two
readers could have been one. On Android they differ in kind: fonts are not files
at all, they live inside the APK and come out through \texttt{AAssetManager}.
Because \texttt{exeDir()} returns \texttt{""} there, the expression
\texttt{exeDir() + "fonts/NewCM10-Regular.otf"} \emph{is} exactly the asset
name. That is what lets an application load its typeface with no \texttt{\#ifdef}
and no knowledge that an APK exists.

Neither reader is for user data. Files the user chose --- documents, music,
images --- are ordinary absolute paths on every platform.

\subsection{Window and monitor}

\begin{longtable}{@{}p{0.42\textwidth}p{0.52\textwidth}@{}}
\toprule
\textbf{Method} & \textbf{Contract} \\
\midrule
\endhead
\texttt{MonitorInfo primaryMonitor() const} &
\texttt{bounds} is the full monitor rectangle; \texttt{workArea} subtracts
taskbars. They are equal everywhere except Windows, because no other platform
here exposes the concept. \\
\addlinespace
\texttt{SafeInsets safeInsets() const} \newline \emph{(default: all zeros)} &
Pixels along each edge the application must not draw in. Zero on both desktops
--- not a stub, a true answer. Re-read every layout pass; a rotation or a fold
changes it without a restart. \\
\addlinespace
\texttt{void adaptToCurrentMonitor()} &
Re-fit if the window moved to another monitor. \textbf{Real no-op on Wayland}:
a client cannot ask which monitor it is on. \\
\addlinespace
\texttt{void snapToEdge(SnapEdge edge)} &
Move the window to \texttt{Left}, \texttt{Right}, \texttt{Top},
\texttt{Bottom} or \texttt{Center} of its monitor. \textbf{Real no-op on
Wayland}: a client cannot position itself. \\
\addlinespace
\texttt{void invalidate()} &
Schedule a repaint. \\
\addlinespace
\texttt{void setCursor(CursorShape)} &
\texttt{Arrow}, \texttt{Hand} or \texttt{Text} --- what a hit-test can honestly
answer. Backends collapse a repeat of the shape already showing, so callers need
not track it. \\
\addlinespace
\texttt{void setKeepAwake(bool)} &
Hold the display awake. Needs a compositor-side protocol on Wayland
(\texttt{zwp\_idle\_inhibit\_manager\_v1}); a no-op where that is absent, since
nothing else can ask. \\
\addlinespace
\texttt{void* secondaryWindowHandle()} \newline \emph{(default:
\texttt{nullptr})} &
Opaque handle so a \emph{second} top-level window can share the platform's
connection rather than opening a second one. \\
\bottomrule
\end{longtable}

\subsection{Time, threads and input registration}

\begin{longtable}{@{}p{0.42\textwidth}p{0.52\textwidth}@{}}
\toprule
\textbf{Method} & \textbf{Contract} \\
\midrule
\endhead
\texttt{void startTimer(int id, int ms)} \newline
\texttt{void stopTimer(int id)} &
A repeating timer, reported as \texttt{onTimer(id)}. See the narrowing note
below. \\
\addlinespace
\texttt{void postAppEvent(int id,} \newline \texttt{\ \ intptr\_t p1 = 0,
intptr\_t p2 = 0)} &
\textbf{Safe to call from any thread.} The three integers come back out of
\texttt{onAppEvent()} unread, on the UI thread, inside \texttt{pump()}. This is
the portable equivalent of posting a private window message to yourself. \\
\addlinespace
\texttt{void registerHotkey(} \newline \texttt{\ \ int id, int keyCode)}
\newline \emph{(default: no-op)} &
Ask for \texttt{Alt+<keyCode>} to arrive as \texttt{onHotkey(id)}. Call
\emph{after} \texttt{init()}: Win32 needs the window to exist. \\
\addlinespace
\texttt{std::string} \newline \texttt{\ \ launchArgument() const}
\newline \emph{(default: empty)} &
What the platform was launched \emph{with}. Empty on both desktops, and empty is
the honest answer there. \\
\bottomrule
\end{longtable}

\paragraph{A documented narrowing.} Only the Win32 backend currently
distinguishes timer ids, because \texttt{SetTimer} requires a real one. Wayland
and Android each have a single timer file descriptor and remember the last id
handed to them rather than honouring it. \textbf{One timer per host is therefore
the supported contract today.} Supporting several needs a descriptor map in
those two backends --- a change to make on the day a second timer exists, not
before.

\subsection{Errors}

\begin{lstlisting}
void showErrorMessage(const std::string& title, const std::string& msg);
\end{lstlisting}

\noindent This is the startup-failure fallback for the case where the application cannot
draw its own error --- Vulkan failed to initialise, a device was not found. It
is a native message box on Windows and a \texttt{stderr} line on Linux.

It is \emph{not} the way to report ordinary failures. Anything the application
can still draw, it should draw.

\subsection{Types}

\begin{lstlisting}
enum class SnapEdge   { Left, Right, Top, Bottom, Center };
enum class CursorShape { Arrow, Hand, Text };

struct MonitorInfo { LayoutRect bounds; LayoutRect workArea; };
struct SafeInsets  { int top = 0, bottom = 0, left = 0, right = 0; };
\end{lstlisting}

\texttt{SnapEdge::Center} is a fifth value rather than a pair of flags because
centring is both axes at once, and a flag pair would admit states that mean
nothing.

%======================================================================
\section{Reference: \texttt{AppView}}

What a platform reports. \textbf{Almost everything has an empty default}, and
that is a design decision with two reasons: an application that takes no
keyboard should not have to write empty overrides to say so, and a backend that
learns to report something new must not break every existing consumer on the day
it does.

Only two are pure virtual, being the two an application cannot meaningfully
lack.

\subsection{Required}

\begin{longtable}{@{}p{0.34\textwidth}p{0.60\textwidth}@{}}
\toprule
\endhead
\texttt{void onHostResized()} &
An explicit size change: relayout, and tell the renderer its surface changed. \\
\addlinespace
\texttt{void shutdown()} &
Teardown before the window and the renderer die. \\
\bottomrule
\end{longtable}

\subsection{Window state}

\begin{longtable}{@{}p{0.34\textwidth}p{0.60\textwidth}@{}}
\toprule
\endhead
\texttt{void onHostLayoutInvalidated()} &
A \emph{routine} resize: relayout, but do not tell the renderer. The
distinction is real on Wayland, where a \texttt{configure} arrives for reasons
that are not a new drawable size. \\
\addlinespace
\texttt{void onHostExposed()} &
Became visible or was uncovered. Nothing changed but the frame is stale. \\
\addlinespace
\texttt{void adaptToCurrentMonitor()} &
The window may have moved to a different monitor. Windows only. \\
\bottomrule
\end{longtable}

\subsection{The surface can die and come back}

\begin{longtable}{@{}p{0.34\textwidth}p{0.60\textwidth}@{}}
\toprule
\endhead
\texttt{void onSurfaceLost()} &
Release every GPU object: swapchain, textures, atlases. \\
\addlinespace
\texttt{bool onSurfaceRecreated()} \newline \emph{(default: true)} &
Rebuild them against the host's \emph{new} surface. Returning false means the
application must not be drawn. \\
\bottomrule
\end{longtable}

\textbf{Neither desktop backend ever calls these} --- they are dead code on
Windows and Linux by construction, because a desktop window is born once and
dies once. On Android the surface dies \emph{every time} the user leaves the
application.

The split they enforce is the one that matters:

\begin{quote}
CPU-side state survives. GPU-side state does not.
\end{quote}

Getting it wrong is invisible until the second visit to the application, when
every string on screen is gone and no error was logged anywhere. A stale
``already uploaded'' flag is the usual culprit.

\subsection{Pointer and keyboard}

\begin{longtable}{@{}p{0.40\textwidth}p{0.54\textwidth}@{}}
\toprule
\endhead
\texttt{void onMouseMove(int x, int y)} & \\
\texttt{void onMouseLeave()} & \\
\texttt{void onLButtonDown(int x, int y)} & \\
\texttt{void onLButtonDblClk(int x, int y)} & \\
\texttt{void onMouseWheel(int x, int y,} \newline \texttt{\ \ int delta)} &
\texttt{delta} is in Win32 \texttt{WHEEL\_DELTA} units (120 per notch) on every
platform, so scroll-step arithmetic is written once. \\
\addlinespace
\texttt{void onNavBack()} \newline \texttt{void onNavForward()} &
The two extra buttons on the side of a mouse
(\texttt{XBUTTON1}/\texttt{XBUTTON2}, \texttt{BTN\_SIDE}/\texttt{BTN\_EXTRA}). \\
\addlinespace
\texttt{void onDragEnd(int dx, int dy)} &
A drag that has \emph{ended}: the pointer travelled \texttt{dx,dy} with the
button held and was released. A tap never reaches this. \\
\addlinespace
\texttt{void onKeyDownPortable(int keyCode)} &
In \texttt{vk\_canvas}'s portable \texttt{key::} space, never a platform's own
keycodes. \\
\addlinespace
\texttt{void onCharPortable(uint32\_t cp)} &
Text entry, one codepoint. \\
\addlinespace
\texttt{void onHotkey(int hotkeyId)} &
A hotkey the application registered. \\
\bottomrule
\end{longtable}

\paragraph{Why \texttt{onDragEnd} is reported rather than reconstructed.} Every
platform already tracks it --- Win32 from \texttt{WM\_LBUTTONDOWN}/\texttt{UP},
Wayland from \texttt{PointerAction::Down}/\texttt{Up}, Android from its own
touch slop. Reconstructing it in the application would mean a fourth
implementation that disagrees with all three. The rule this draws:

\begin{quote}
The \emph{slop} is a host property. The \emph{meaning} is an application
property.
\end{quote}

On a touch screen the host also synthesises hover from the finger's position,
which is the only honest hover a touch screen has.

\subsection{The carried identifiers}

\begin{longtable}{@{}p{0.40\textwidth}p{0.54\textwidth}@{}}
\toprule
\endhead
\texttt{void onTimer(int timerId)} &
A timer fired. \\
\addlinespace
\texttt{void onAppEvent(int id,} \newline \texttt{\ \ intptr\_t p1, intptr\_t
p2)} &
A cross-thread completion posted via \texttt{postAppEvent()}, delivered on the
UI thread from inside \texttt{pump()}. \\
\addlinespace
\texttt{void onHostReady()} &
The first moment the application is fully constructed \emph{and} the host is
running: after \texttt{create()} returned, from inside the first
\texttt{pump()}. \\
\bottomrule
\end{longtable}

\paragraph{\texttt{onHostReady()} may fire more than once.} Platforms that
suspend and resume call it again, so an application doing one-time work there
must guard it \emph{itself} --- preferably on the state it is about to create
rather than on a boolean, which is the difference between a guard and a flag
that can go stale.

Its purpose is the launch argument. A typical implementation:

\begin{lstlisting}
void MyApp::onHostReady() {
    if (alreadyConfigured()) return;
    const std::string arg = host_->launchArgument();
    if (!arg.empty()) configureFrom(arg);
}
\end{lstlisting}

%======================================================================
\section{Reference: the portable modules}

Four small modules that every application of this shape needs, and that are
almost always written badly the first time.

\subsection{\texttt{app\_paths} --- where files live}

\begin{lstlisting}
namespace app_paths {
    const std::string& exeDir();    // read-only, shipped
    const std::string& stateDir();  // everything WRITTEN
}
\end{lstlisting}

Both are UTF-8 with a trailing separator. The rule is the split:

\begin{description}[leftmargin=1.5em]
\item[\texttt{exeDir()}] Read-only data shipped beside the binary: fonts,
  shaders, default configuration. Nothing is ever written here, which is what
  lets a package install it root-owned.
\item[\texttt{stateDir()}] Everything the application writes: databases, logs,
  caches. Created on first call.
\end{description}

\textbf{They are the same directory by default}, so a build tree and a portable
archive stay one self-contained folder you can move anywhere. Defining
\texttt{APP\_SHELL\_STATE\_HOME} at build time moves the writable half to
\texttt{\$HOME/<that name>/}, which is what a system package needs, because
\texttt{/opt} is root-owned and a cache is per-user regardless.

Deliberately \emph{not} XDG. A plain dotdir in the \texttt{\textasciitilde/.ssh}
tradition, depending on no specification.

If the home directory is unknown or the directory cannot be created and written
to, \texttt{stateDir()} falls back to \texttt{exeDir()} --- a stripped
environment degrades to the old behaviour instead of failing to start.

On Android neither question means anything: the assets are in the APK and the
writable directory is handed to the application for free. There a different
implementation answers the same two functions, and the desktop one is not
compiled at all. Building both is a duplicate-symbol link error, which is
exactly the right failure for two implementations of one seam.

\subsection{\texttt{ui\_metrics} --- one scale, five type roles}

\begin{lstlisting}
struct UiTextSizes { float caption, secondary, body, title, header; };

struct UiMetrics {
    float       scale;   // 1.0 at or below the reference, proportional above
    UiTextSizes text;
    float space(float authored) const;   // pads, gaps, row heights
    float stroke(float authored) const;  // snapped to whole device pixels
};

UiMetrics computeUiMetrics(float contentShortSide);
\end{lstlisting}

Two decisions here are load-bearing.

\paragraph{It takes the window's SHORT SIDE, not its height.} These are the same
number while every window is wider than tall, and stop being the same the moment
a vertical layout exists. A $1080\times1920$ monitor would otherwise scale
$1.78\times$ and a $1080\times2400$ phone $2.22\times$ --- inflating every bar
and every type role because the screen is \emph{tall}, not because it is
\emph{big}. What limits how much fits is the short side.

\paragraph{Every role derives from one clamped value.} The smallest role is
pinned to the fonts' own geometric legibility floor and each larger role is that
floor times a fixed ratio. Because all five come from one number, the scale
clamps \emph{uniformly} below the reference height instead of the smallest role
clamping alone and squashing the hierarchy flat. That flattening is the specific
defect this module was written to replace: seven independently hand-tuned
percentages, every one of which sat below the legibility floor at its own
reference size.

\texttt{stroke()} snaps to whole device pixels, so a one-pixel rule never lands
blurred across two --- continuous scaling would put it at $1.33$\,px on a 1440p
screen.

\paragraph{The generated header.} \texttt{ui\_metrics.hh} includes
\texttt{ui\_min\_text\_size.gen.h}, produced at build time from \emph{the
consuming application's own fonts}: the emitted floor is the worst case across
them. It is therefore a per-application number and not a library constant. See
Section~\ref{sec:mintext}.

\subsection{\texttt{ui\_orientation} --- which way round the window is}

\begin{lstlisting}
enum class UiOrientation { Horizontal, Vertical };

UiOrientation autoOrientationFor(int windowW, int windowH);

struct UiOrientationState {
    bool          autoEnabled = true;
    UiOrientation manual      = UiOrientation::Horizontal;
    UiOrientation resolve(int windowW, int windowH) const;
    void toggleManual(int windowW, int windowH);
};
\end{lstlisting}

\textbf{Derived from the shape the window already has, never asked of the OS.}
That single property is why it behaves identically on Wayland (where a client
cannot size or position itself), on Windows, and on Android (where a device
rotation arrives as an ordinary resize).

It replaced a pair of ``modes'' that were really window rectangles wearing a
mode's name, and whose Wayland implementation could only ask the compositor for
a size and wait for an asynchronous \texttt{configure} to come back --- which is
why toggling the layout used to feel slow.

A tie ($w = h$) resolves to \texttt{Horizontal} deliberately: something has to
win. Degenerate sizes --- a Wayland surface configured $0\times0$ before its
first real configure --- fall out of the same comparison rather than needing a
guard, so the answer is defined rather than accidental.

\texttt{toggleManual()} seeds \texttt{manual} from what automatic \emph{would}
have said before flipping it. Without that seeding, the first key press on a
vertical monitor appears to do nothing.

\subsection{\texttt{layout\_rect} and \texttt{color}}

Two POD types that exist to delete two Windows dependencies:

\begin{lstlisting}
struct LayoutRect { int left = 0, top = 0, right = 0, bottom = 0; };

using ColorRef = uint32_t;                       // 0x00BBGGRR, as COLORREF was
constexpr ColorRef RGB(uint8_t r, uint8_t g, uint8_t b);
constexpr uint8_t  GetRValue(ColorRef c);        // and G, B
\end{lstlisting}

Both keep Win32's member names and bit layout on purpose, so porting existing
code touches no call site.

\texttt{color.hh} \texttt{\#undef}s \texttt{RGB}, \texttt{GetRValue},
\texttt{GetGValue} and \texttt{GetBValue} before declaring its own. Those are
\texttt{wingdi.h} \emph{textual macros}: if \texttt{<windows.h>} reached the
translation unit first, they do not shadow the declarations below, they break
them outright by substitution.

%======================================================================
\section{Platform notes}

\subsection{Win32 (\texttt{os/win32\_host.cc})}

A real \texttt{HWND} with a classic \texttt{WndProc}. Owns \texttt{WinMain},
which raises DPI awareness, opens the log, installs a minidump crash handler,
raises the system timer resolution to 1\,ms, initialises COM, and then calls
\texttt{app\_shell\_main()}.

It is the only backend that distinguishes timer ids, and the only one where
\texttt{registerHotkey()} is genuinely \emph{system-wide}: a registered hotkey
fires with the application in the background. Hotkeys are unregistered on
teardown, tracked by the ids that registered successfully rather than by a
fixed list.

\texttt{Host} carries one Windows-only method,
\texttt{HWND~nativeHandle()~const}, behind \texttt{\#ifdef~\_WIN32}. It exists because mouse-leave detection
(\texttt{TrackMouseEvent}) needs a real handle and there is no portable
equivalent. It is not part of the portable surface and nothing portable may call
it.

\subsection{Wayland (\texttt{os/wayland\_host.cc})}

Built on \texttt{vk\_canvas}'s \texttt{WaylandDisplay}/\texttt{WaylandWindow}. A
\texttt{timerfd} for the timer and an \texttt{eventfd}-woken queue for
cross-thread events, both drained from \texttt{pump()}. Owns \texttt{main()},
which opens the log and installs signal handlers before calling
\texttt{app\_shell\_main()}.

What Wayland genuinely cannot do, and which is therefore answered honestly
rather than faked:

\begin{itemize}[leftmargin=*]
\item \texttt{snapToEdge()} --- a client cannot position itself.
\item \texttt{adaptToCurrentMonitor()} --- a client cannot ask which monitor it
  is on.
\item \texttt{primaryMonitor().workArea} --- no taskbar concept is exposed, so
  it equals \texttt{bounds}.
\item \texttt{registerHotkey()} --- no cross-compositor global-hotkey protocol
  exists, so it becomes a focused-window-only check against the registered
  table. \textbf{A documented narrowing, not a silent drop.}
\end{itemize}

These are no-ops with comments, not bugs. The difference matters: a reader who
finds a silent no-op assumes it is broken.

\paragraph{One logging detail worth keeping.} \texttt{openLogFile()}
\texttt{dup2}s \texttt{stderr} onto \texttt{stdout}'s descriptor. Opening the
same path twice gives the two streams independent offsets, and one silently
erases the other's writes.

\subsection{Android (\texttt{os/android\_host.cc})}
\label{sec:android}

Modelled directly on the Wayland host: bionic has \texttt{timerfd} and
\texttt{eventfd}, so the timer and the event queue are the same calls, merely
registered with \texttt{ALooper\_addFd()} instead of a hand-rolled
\texttt{poll()}.

The application supplies its own \texttt{android\_main()}:

\begin{lstlisting}
void android_main(android_app* state) {
    MyApp app;
    // "my_arg" is the intent extra launchArgument() should return; the
    // second string is the fallback when it is absent or empty.
    if (!app.create(std::make_unique<AndroidHost>(state, "my_arg", "/default")))
        return;
    app.run();
    app.shutdown();
}
\end{lstlisting}

\paragraph{Be loud here.} When \texttt{android\_main()} returns, the activity
finishes --- so every startup failure looks identical from outside: the
application opens and closes. There is no window left to put an error in and no
console to print one to. Log the phases.

\paragraph{Standard output goes to \texttt{/dev/null} on Android.}
\texttt{AndroidHost} therefore \texttt{dup2}s \texttt{stdout} and
\texttt{stderr} onto a pipe whose reader thread forwards whole lines to logcat
under the tag \texttt{APP\_SHELL\_APP\_NAME}. Without it, every
\texttt{printf} in the application and in every library it links is silently
discarded --- on the one platform with no terminal. Partial lines are held
rather than flushed, so messages are never torn at a read boundary.

\paragraph{Touch is translated, not pretended.} Under a 24\,px slop the gesture
is a tap and the press is delivered at \emph{release}, so a finger that slides
off cancels, the way a button works everywhere. Past the slop the gesture
becomes wheel deltas for good and the press never happens.

\paragraph{Bundled JNI plumbing.} Four modules travel with this host:
\texttt{launch\_intent} (read a string extra off the launch intent),
\texttt{safe\_area} (the display cutout), \texttt{storage\_permission}
(all-files access, asked at most once), and \texttt{app\_paths\_android}. The permission is asked once per process
\emph{and} only when the system says it is not already granted; both halves are
needed, or a user who declines is re-asked on every resume.

\subsection{A fourth host: headless}

Nothing prevents an application from implementing \texttt{Host} itself and
passing it in. The reference consumer does exactly this for deterministic
screenshot testing: a host backed by \texttt{VK\_EXT\_headless\_surface} that
answers ``there is no window'' to everything else, so the entire real
application --- Vulkan initialisation, database restore, font baking, layout ---
runs against a surface that presents to nothing.

This is the strongest argument for the seam being drawn in the right place. If
\texttt{Host} were not a genuine interface, that tool would have to
re-implement the application, and would stop describing it the moment it moved.

%======================================================================
\section{Build integration}

\subsection{Targets}

\begin{center}
\begin{tabular}{@{}ll@{}}
\toprule
Target & Contents \\
\midrule
\texttt{app\_shell}         & portable: \texttt{app\_paths}, \texttt{ui\_metrics}, \texttt{ui\_orientation}, headers \\
\texttt{app\_shell\_win32}   & the Win32 host \\
\texttt{app\_shell\_wayland} & the Wayland host \\
\texttt{app\_shell\_android} & the Android host and its JNI plumbing \\
\bottomrule
\end{tabular}
\end{center}

A consumer links \texttt{app\_shell} plus \emph{exactly one} of the others. Only
the target for the platform being configured is defined, so linking the wrong
one is a configure error rather than a link error.

\subsection{Cache variables}

Set them \emph{before} \texttt{add\_subdirectory()}. All have working defaults,
so the library configures and builds on its own.

\begin{center}
\begin{tabular}{@{}lll@{}}
\toprule
Variable & Default & Effect \\
\midrule
\texttt{APP\_SHELL\_APP\_NAME}   & \texttt{app}          & log file base name; Android logcat tag \\
\texttt{APP\_SHELL\_WIN\_CLASS}  & \texttt{AppShellMain} & Win32 window class \\
\texttt{APP\_SHELL\_STATE\_HOME} & (empty)               & dir under \texttt{\$HOME} for written files \\
\texttt{APP\_SHELL\_NATIVE\_APP\_GLUE\_DIR} & --- & required on Android \\
\bottomrule
\end{tabular}
\end{center}

\texttt{APP\_SHELL\_NATIVE\_APP\_GLUE\_DIR} has no default and no sensible
guess: \texttt{native\_app\_glue} ships inside the NDK, so only the consumer
knows where it is. Omitting it is a \texttt{FATAL\_ERROR} with that sentence in
it, rather than a confusing missing-header failure five files later.

\subsection{The text-size floor generator}
\label{sec:mintext}

\begin{lstlisting}[language=]
include(AppShellMinTextSize)
app_shell_generate_min_text_size(
    ${CMAKE_SOURCE_DIR}/assets/fonts/MyFace-Regular.otf
    ${CMAKE_SOURCE_DIR}/assets/fonts/MyFace-Bold.otf)
add_dependencies(app_shell generate_ui_min_text_size)
\end{lstlisting}

This emits \texttt{ui\_min\_text\_size.gen.h} into the build tree's
\texttt{generated/} directory, which \texttt{ui\_metrics.hh}
includes. The value is derived from the listed
faces' own outline geometry: the smallest device-pixel size at which their
strokes are still above what antialiased rendering can represent.

\textbf{List exactly the faces the application loads.} The emitted floor is the
worst case across them, so one left out can render below its own floor.

Two constraints:

\begin{itemize}[leftmargin=*]
\item \textbf{It cannot run under a cross-compiler}, because the tool must
  execute on the build machine. An Android build copies the header out of a
  desktop build tree instead, which means a desktop build must have happened
  first.
\item The dependency must be stated by the consumer --- \\
  \texttt{add\_dependencies(app\_shell generate\_ui\_min\_text\_size)} --- \\
  because the target does not exist yet when CMake reads this library's own
  \texttt{CMakeLists.txt}. The consumer owns the font list, so the consumer owns
  the ordering.
\end{itemize}

%======================================================================
\section{Testing}

Two assert-based executables, Debug-only, no framework:

\begin{lstlisting}[language=]
./build/<tree>/app_shell_build/ui_metrics_test
./build/<tree>/app_shell_build/ui_orientation_test
\end{lstlisting}

They \emph{compile} the sources rather than linking \texttt{app\_shell}, which
is the property that keeps them honest: a test can never quietly start depending
on Vulkan. Test sources \texttt{\#undef NDEBUG} so asserts survive an optimised
build.

The hosts have no automated tests and are unlikely ever to have useful ones ---
they are thin adapters over three operating systems. What \emph{is} testable, and
what the reference consumer does, is the whole application through an injected
headless host, compared image by image against a known-good set. That catches
regressions in this library indirectly but very effectively: a change here that
alters layout by one pixel shows up immediately.

%======================================================================
\section{Known limitations}

Stated rather than hidden, in rough order of how likely they are to matter.

\begin{enumerate}[leftmargin=*]
\item \textbf{One timer per host.} Two of the three backends ignore the timer
  id. See Section~4.4.
\item \textbf{No IME on Android.} \texttt{onCharPortable()} is never fed there,
  so a text field can be seen and not typed into.
\item \textbf{Hotkeys are Alt-only.} \texttt{registerHotkey()} takes a key code
  and implies Alt. A modifier mask is the obvious extension.
\item \textbf{No run loop is provided.} Deliberate today --- what counts as work
  is the application's question --- but a default loop for applications that do
  not care would be a reasonable addition.
\item \textbf{No multi-window support.} \texttt{secondaryWindowHandle()} is an
  escape hatch for an application that wants a second top-level window, not an
  API for one. There is exactly one \texttt{Host} per process.
\item \textbf{No clipboard, no drag-and-drop, no file dialogs.} An application
  that needs them either implements them itself or extends this library.
\end{enumerate}

%======================================================================
\section{History, and why the shapes are what they are}

This library was extracted from a working application rather than designed in
advance, and every unusual decision in it is a scar. Recording them is cheaper
than rediscovering them.

The extraction was done in two deliberate stages so that any regression would be
bisectable: first \emph{generalise in place}, changing types and call directions
while every file stayed where it was; then \emph{move}, changing no logic at all.

Four couplings had to be cut:

\begin{enumerate}[leftmargin=*]
\item \texttt{Host::init()} took the concrete application class. Now it takes
  \texttt{AppView*}.
\item Cross-thread events and timer ids were enumerations declared in
  \texttt{host.hh}, and every backend switched on them --- the same dispatch
  written out four times. Both are plain \texttt{int} now, and the switch exists
  once, in the application.
\item \texttt{snapToEdge()} took the application's hotkey id, so the Win32 host
  held a table mapping one application's keys to window geometry. It takes a
  \texttt{SnapEdge}; hotkeys are registered by the application.
\item The Android host read a named intent extra and called the application's
  own domain method with it. It is told the key now, and states the string.
\end{enumerate}

\paragraph{How the result was verified.} The reference application renders
twenty distinct UI states through a headless host. Before any change, those
twenty images were captured as a baseline; after each stage, and again after a
clean rebuild, they were compared \emph{byte for byte} and were identical. That,
plus a suite of pure-logic tests and a bit-exactness gate on the audio path, is
the evidence that lifting this code changed nothing about the program it was
lifted out of.

\vfill
\begin{center}
\rule{0.4\textwidth}{0.4pt}
\end{center}

\end{document}
